\section{Experiments}
\label{sec:experiments}

\subsection{Experiment Setup}

\textbf{Models and Datasets.} To comprehensively evaluate the universality of our algorithm, we selected the LLaMA-3~\cite{llama3} series and Qwen3~\cite{qwen3} series as foundation models, covering parameter scales from 8B to 70B and including both Dense and Mixture-of-Experts (MoE) architectures. Beyond standard perplexity evaluation on Wikitext2~\cite{wikitext2}, we assessed TORQ on a suite of zero-shot tasks, including ARC-Challenge and ARC-Easy~\cite{arc}, HellaSwag~\cite{hellaswag}, LAMBADA~\cite{lambada}, PIQA~\cite{piqa}, and WinoGrande~\cite{winogrande}. Furthermore, we evaluated performance on more challenging reasoning benchmarks, such as the multi-disciplinary knowledge task MMLU~\cite{mmlu}, the mathematical reasoning benchmark GSM8K~\cite{gsm8k}, and the code generation benchmark HumanEval~\cite{humaneval}.

\textbf{Baseline Methods.} We compared TORQ with current state-of-the-art low-bit quantization methods, including RTN, GPTQ~\cite{gptq}, QuaRot~\cite{quarot}, OSTQuant~\cite{ostquant}, FlatQuant~\cite{flatquant}, and BRQ~\cite{brq}. For fair comparison, all baseline methods were evaluated without GPTQ post-processing, examining solely the quantization performance of the methods themselves.

\textbf{Implementation Details.} All experiments were conducted on NVIDIA RTX 5090 GPUs. We applied MXFP4 quantization to both activations and weights. For calibration, we sampled 128 text segments from the Wikitext2 dataset. As TORQ is an efficient Post-Training Quantization (PTQ) framework, no fine-tuning was performed. For dense models (LLaMA3 and Qwen3), we evaluated MMLU, GSM8K, and HumanEval alongside other zero-shot tasks using the open-source \texttt{lm-evaluation-harness}~\cite{lm_eval}.

\subsection{Main Results}

Table \ref{tab:main_results} presents the main experimental results on the LLaMA3 and Qwen3 model families. \textbf{TORQ significantly outperforms SOTA methods across all tested models and tasks.} 
For example, on LLaMA3-8B, RTN and GPTQ nearly failed in the MXFP4 format (PPL > 280), confirming the challenge of directly applying MXFP4. In contrast, TORQ reduced the WikiText2 perplexity to \textbf{8.57}, surpassing the runner-up methods FlatQuant (8.92) and OSTQuant (9.29). In terms of average accuracy, TORQ achieved \textbf{69.32\%} on LLaMA3-8B, improving by over 37 percentage points compared to the baseline (RTN) and outperforming FlatQuant (68.34\%).

\begin{table*}[h]
\centering
\caption{Main results on LLaMA3 and Qwen3 families. "Wiki2" reports Perplexity ($\downarrow$), and "Avg Acc" reports average accuracy on zero-shot tasks ($\uparrow$).}
\label{tab:main_results}
\resizebox{\textwidth}{!}{
\begin{tabular}{lcccccccccc}
\toprule
 & \multicolumn{2}{c}{\textbf{LLaMA3-8B}} & \multicolumn{2}{c}{\textbf{LLaMA3-70B}} & \multicolumn{2}{c}{\textbf{Qwen3-8B}} & \multicolumn{2}{c}{\textbf{Qwen3-14B}} & \multicolumn{2}{c}{\textbf{Qwen3-32B}} \\
\cmidrule(lr){2-3} \cmidrule(lr){4-5} \cmidrule(lr){6-7} \cmidrule(lr){8-9} \cmidrule(lr){10-11}
\textbf{Method} & Wiki2 & Avg Acc & Wiki2 & Avg Acc & Wiki2 & Avg Acc & Wiki2 & Avg Acc & Wiki2 & Avg Acc \\
\midrule
Baseline (FP16) & 6.14 & 73.18 & 2.87 & 79.95 & 9.00 & 70.43 & 8.64 & 73.81 & 7.61 & 74.82 \\
\midrule
RTN & 354.37 & 32.03 & 3180.43 & 30.92 & 3750.93 & 30.73 & 2219.90 & 35.22 & 274.93 & 38.40 \\
GPTQ & 282.29 & 33.60 & 1532.08 & 34.82 & 1633.38 & 35.54 & 368.24 & 39.91 & 59.31 & 57.62 \\
QuaRot & 10.16 & 62.13 & 15.38 & 62.87 & 23.50 & 55.35 & 20.11 & 60.21 & 18.86 & 61.23 \\
OSTQuant & 9.29 & 65.92 & 6.01 & 73.93 & 19.75 & 54.40 & 16.85 & 64.06 & 14.62 & 67.19 \\
FlatQuant & 8.92 & 68.34 & OOM & OOM & 11.21 & 66.96 & 13.32 & 69.87 & OOM & OOM \\
BRQ & 9.14 & 67.76 & - & - & 13.70 & 64.73 & - & - & - & - \\
\textbf{TORQ (ours)} & \textbf{8.57} & \textbf{69.32} & \textbf{3.49} & \textbf{76.28} & \textbf{10.26} & \textbf{67.75} & \textbf{12.87} & \textbf{70.51} & \textbf{8.43} & \textbf{73.63} \\
\bottomrule
\end{tabular}
}
\end{table*}

\subsection{Ablation Studies}

\textbf{Component Analysis.} To verify the contribution of individual components, we conducted ablation studies on Qwen3-8B (Table \ref{tab:ablation_component}). 
\begin{itemize}
    \item Using only Inter-block Rotation ($\mathbf{R}_{\text{inter}}$) reduced PPL from 3750.93 (RTN) to 24.73, indicating that resolving inter-block variance imbalance is fundamental for recovering model functionality.
    \item Using only Intra-block Rotation ($\mathbf{R}_{\text{intra}}$) lowered PPL to 20.36, demonstrating significant gains from optimizing codebook utilization.
    \item The complete TORQ ($\mathbf{R}_{\text{inter}} + \mathbf{R}_{\text{intra}}$) further reduced PPL to 10.26, boosting GSM8K accuracy to 85.95\%. This proves that macro-variance equilibrium and micro-geometric alignment are orthogonal and complementary.
\end{itemize}

\begin{table*}[h]
\centering
\caption{Ablation study of rotation components on Qwen3-8B.}
\label{tab:ablation_component}
{
\begin{tabular}{cc|cccccc}
\toprule
$\mathbf{R}_{\text{inter}}$ & $\mathbf{R}_{\text{intra}}$ & \textbf{Wiki2} & \textbf{ARC-C} & \textbf{ARC-E} & \textbf{Hella} & \textbf{PIQA} & \textbf{GSM8K} \\
\midrule
\multicolumn{2}{c|}{Baseline (FP16)} & 9.00 & 56.74 & 80.93 & 74.98 & 77.37 & 88.25 \\
\midrule
$\times$ & $\times$ & 3750.93 & 22.94 & 25.07 & 28.73 & 52.11 & 48.93 \\
$\checkmark$ & $\times$ & 24.73 & 40.33 & 60.18 & 55.30 & 61.46 & 65.98 \\
$\times$ & $\checkmark$ & 20.36 & 48.41 & 69.17 & 61.01 & 68.83 & 70.84 \\
$\checkmark$ & $\checkmark$ & \textbf{10.26} & \textbf{53.65} & \textbf{79.36} & \textbf{72.06} & \textbf{74.11} & \textbf{85.95} \\
\bottomrule
\end{tabular}
}
\end{table*}

\textbf{Calibration Dataset Robustness.} We evaluated the sensitivity of TORQ to different calibration datasets on Qwen3-14B (Table \ref{tab:ablation_data}). Results show minimal performance fluctuation (PPL variance < 0.8) across WikiText2, C4, and RedPajama, indicating that TORQ is robust and can capture statistical features from small amounts of generic text.

\begin{table}[h]
\centering
\caption{Robustness across different calibration datasets on Qwen3-14B.}
\label{tab:ablation_data}
\resizebox{0.5\textwidth}{!}{
\begin{tabular}{lcccc}
\toprule
\textbf{Calibration Data} & \textbf{Wiki2} & \textbf{ARC-C} & \textbf{ARC-E} & \textbf{Hella} \\
\midrule
Wiki2 & 12.87 & 55.27 & 81.74 & 76.83 \\
C4 & 13.63 & 56.11 & 81.09 & 78.02 \\
Red\_Pajama & 13.11 & 55.78 & 80.93 & 77.37 \\
\bottomrule
\end{tabular}
}
\end{table}

\textbf{Performance on Hard Tasks.} We assessed "deep intelligence" on MMLU, GSM8K, and HumanEval (Table \ref{tab:hard_tasks}). Traditional methods suffered severe degradation; for instance, GPTQ on Qwen3-8B achieved only 67.37\% on GSM8K and 30.84\% on HumanEval. TORQ recovered GSM8K accuracy to \textbf{86.82\%} (vs. 88.25\% FP16) and HumanEval to \textbf{61.15\%} (vs. 64.63\% FP16), thanks to the micro-alignment rotation preserving fine-grained features critical for reasoning.


\begin{table}[h]
\centering
\caption{Performance on complex reasoning tasks (Qwen3-8B).}
\label{tab:hard_tasks}
\resizebox{0.5\textwidth}{!}{
\begin{tabular}{lcccc}
\toprule
\textbf{Method} & \textbf{MMLU} & \textbf{GSM8K} & \textbf{HumanEval} & \textbf{MathQA} \\
\midrule
Baseline & 74.71 & 88.25 & 64.63 & 49.65 \\
GPTQ & 50.29 & 67.37 & 30.84 & 36.92 \\
QuaRot & 68.88 & 79.30 & 53.37 & 41.12 \\
\textbf{TORQ (ours)} & \textbf{70.82} & \textbf{86.82} & \textbf{61.15} & \textbf{47.94} \\
\bottomrule
\end{tabular}
}
\end{table}
\vspace{-3mm}

\textbf{Generalizability to Other Formats.} We applied TORQ to MXINT4 and NVFP4 formats (Table \ref{tab:formats}). TORQ yielded significant improvements regardless of integer (MXINT4) or floating-point (NVFP4) constraints. Notably, on NVFP4, TORQ improved GSM8K accuracy on Qwen3-8B from 42.67\% (RTN) to 87.08\%, confirming the format-agnostic nature of our geometric optimization.

\begin{table}[!htbp]  % 强制优先当前位置，表格不跑飞
\centering
\renewcommand\arraystretch{1.15}  % 行间距不变，保留你的优化
\caption{Applicability to different microscaling formats (Qwen3-8B).}
\label{tab:formats}  % label在caption后，无hyperref锚点警告
\resizebox{0.5\textwidth}{!}{
\begin{tabular}{c c c cccc}  % ✅ 核心修复：列数精准匹配，无多余&符号
\toprule
\textbf{Format} & \textbf{Method} & \textbf{ARC-C} & \textbf{GSM8K} & \textbf{Hella} & \textbf{PIQA} & \textbf{AVG} \\
\midrule
\multirow{2}{*}{MXFP4} & RTN & 22.94 & 35.87 & 28.73 & 52.11 & 34.92 \\
                       & \textbf{TORQ} & \textbf{53.65} & \textbf{86.82} & \textbf{72.06} & \textbf{74.11} & \textbf{71.66} \\
\midrule
\multirow{2}{*}{MXINT4} & RTN & 23.83 & 37.90 & 30.15 & 50.83 & 35.68 \\
                        & \textbf{TORQ} & \textbf{51.98} & \textbf{85.12} & \textbf{72.72} & \textbf{73.51} & \textbf{70.83} \\
\midrule
\multirow{2}{*}{NVFP4} & RTN & 30.17 & 42.67 & 37.33 & 60.12 & 42.57 \\
                       & \textbf{TORQ} & \textbf{54.38} & \textbf{87.08} & \textbf{73.10} & \textbf{75.83} & \textbf{72.60} \\
\bottomrule
\end{tabular}
}
\end{table}

\textbf{Adaptability to MoE Architectures.} Experiments on Qwen3-30B-A3B (MoE) demonstrated TORQ's effectiveness (Table \ref{tab:moe}). On GSM8K, TORQ achieved 83.83\%, far surpassing GPTQ (37.89\%) and QuaRot (74.18\%), proving its ability to mitigate the heterogeneity of activation distributions across experts.

\begin{table*}[h]
\centering
\caption{Performance on MoE Architecture (Qwen3-30B-A3B).}
\label{tab:moe}
% \resizebox{0.9\textwidth}{!}{
\begin{tabular}{lcccccc}
\toprule
\textbf{Method} & \textbf{ARC-E} & \textbf{ARC-C} & \textbf{Hella} & \textbf{GSM8K} & \textbf{WinG} & \textbf{PIQA} \\
\midrule
Baseline & 79.25 & 56.40 & 79.60 & 85.44 & 70.32 & 80.30 \\
RTN & 24.58 & 26.80 & 25.52 & 22.45 & 52.87 & 51.04 \\
GPTQ & 31.15 & 34.29 & 40.19 & 37.89 & 58.90 & 62.81 \\
QuaRot & 70.45 & 53.19 & 71.14 & 74.18 & 62.51 & 71.11 \\
OSTQuant & 75.84 & 54.03 & 77.23 & 80.28 & 65.80 & 75.09 \\
\textbf{TORQ (ours)} & \textbf{77.10} & \textbf{55.81} & \textbf{78.12} & \textbf{83.83} & \textbf{68.36} & \textbf{77.93} \\
\bottomrule
\end{tabular}
% }
\end{table*}

\subsection{Efficiency Analysis}

\textbf{Speedup and Memory Saving.} As shown in Table \ref{tab:speedup}, on Qwen3-30B-A3B with a 2k context length, TORQ (W4A4) achieves a decoding speed of 220 Toks/s, representing a \textbf{3.5$\times$ speedup} over BF16 (63 Toks/s), while reducing memory usage from 56.9GB to 15.3GB.

\begin{table}[h]
\centering
\caption{Inference speedup and memory footprint on Qwen3-30B-A3B.}
\label{tab:speedup}
\begin{tabular}{lcc}
\toprule
\textbf{Precision} & \textbf{Memory (GB)} & \textbf{Decoder Toks/s} \\
\midrule
BF16 & 56.9 & 63 \\
\textbf{TORQ (W4A4)} & \textbf{15.3} & \textbf{220} \\
\bottomrule
\end{tabular}
\end{table}

\textbf{Quantization Time Comparison.} Unlike training-dependent methods, TORQ's analytical construction is extremely efficient. As detailed in Table \ref{tab:time_cost}, TORQ calibrates Qwen3-32B in just 1920 seconds, which is approximately \textbf{10$\times$ faster than FlatQuant} and \textbf{7$\times$ faster than OSTQuant}, highlighting its practical deployment value.

\begin{table}[h]
\centering
\caption{Comparison of quantization calibration time (seconds).}
\label{tab:time_cost}
\begin{tabular}{lccc}
\toprule
\textbf{Model} & \textbf{OSTQuant} & \textbf{FlatQuant} & \textbf{TORQ} \\
\midrule
Qwen3-8B & 2786s & 3247s & \textbf{427s} \\
Qwen3-14B & 5131s & 6489s & \textbf{893s} \\
Qwen3-32B & 14834s & 18619s & \textbf{1920s} \\
\bottomrule
\end{tabular}
\end{table}
